\setchapterimage{figure/Chapter8封面图.jpg}
\setchapterintro{本章利用相位角讨论二阶线性微分方程解的零点，证明 Sturm 分离定理与比较定理，并研究两点边值问题的特征值和特征函数。}
\chapter{Sturm 零点理论与边值问题}\label{chap:sturm}

第三章已由 Wronski 行列式说明两个解是否独立，第五章又给出解对参数的可微依赖性。现在将这两个结果用于二阶方程：先数零点，再通过参数改变研究边值问题。相位角沿用第六章相平面分析的几何思想。

\section{零点与相位角}
设 $q$ 在区间 $I$ 连续，考虑
\[
 x''+q(t)x=0,\qquad y=x'.
\]
除恒零解外，状态 $(x(t),y(t))$ 从不等于零，否则由线性唯一性解恒为零。
在时间区间上可连续选择角度 $\theta$，使
$x=\rho\sin\theta,y=\rho\cos\theta$，$\rho>0$。
\begin{proposition}\label{prop:arnold-angle}
\[
 \theta'=\frac{y^2+q(t)x^2}{x^2+y^2}
 =\cos^2\theta+q(t)\sin^2\theta.
\]
非零解的零点均为单零点且在 $I$ 内无聚点；经过零点时 $\theta$ 以速度 1 穿过某个 $k\pi$。
两个相邻零点之间角度增加恰好 $\pi$。
\end{proposition}
\begin{proof}
在 $y\ne0$ 的小区间，$\tan\theta=x/y$，所以
\[
 \sec^2\theta\,\theta'
 =\frac{x'y-xy'}{y^2}=\frac{y^2+q(t)x^2}{y^2}.
\]
利用 $\sec^2\theta=(x^2+y^2)/y^2$，得到所需公式。
在 $y=0$ 处可改用 $\cot\theta=y/x$，因为此时 $x\ne0$，得到同样公式。
所以公式在整条非零状态轨道上成立。

若 $x(t_0)=0$，初值唯一性禁止 $y(t_0)=0$，故 $x'(t_0)\ne0$，这是单零点。
公式在这里给出 $\theta'(t_0)=1$，所以经过零点时角度总沿增大方向穿过 $k\pi$。
若零点列 $t_j$ 在区间内部聚到 $t_*$，连续性给出 $x(t_*)=0$；
在相邻零点间用 Rolle 定理并取极限，又得 $x'(t_*)=0$，与非零解的唯一性矛盾。

设 $a<b$ 为相邻零点，取 $\theta(a)=k\pi$。
在 $a$ 后角度大于 $k\pi$，且中间不能再取任何整数倍 $\pi$，否则会多出零点。
到 $b$ 前角度必须从下方趋近下一个整数倍，因为在 $b$ 的穿越方向也为增大。
因此 $\theta(b)=(k+1)\pi$，总角增正好为 $\pi$。
这不要求角度在两个零点之间处处单调。
\end{proof}
零点交错的结论需要区分线性相关和线性无关两种情形：
相关解的零点相同，严格交错只对线性无关的解成立。

\begin{theorem}[Sturm 分离定理]
若 $x_1,x_2$ 为线性无关的非零解，则 $x_1$ 的任意两个相邻零点之间
恰有 $x_2$ 的一个零点。
\end{theorem}
\begin{proof}
记 $W=x_1x_2'-x_1'x_2$。由 Liouville 公式，$W$ 为常数；
线性无关性又保证 $W\ne0$。
因此两个状态向量在任意同一时刻都不共线，
对应连续角度之差永远不是 $\pi$ 的整数倍。

设 $t_1<t_2$ 是 $x_1$ 的相邻零点。改变解的符号和角度分支后，可令
$\theta_1(t_1)=0,\theta_1(t_2)=\pi$，并令 $0<\theta_2(t_1)<\pi$。
角度之差连续且不能经过 $0$ 或 $\pi$，所以在整个 $[t_1,t_2]$ 上
$0<\theta_2-\theta_1<\pi$。
特别地，$\theta_2(t_1)<\pi$ 而 $\theta_2(t_2)>\pi$。
介值定理给出 $s\in(t_1,t_2)$ 使 $\theta_2(s)=\pi$，即 $x_2(s)=0$。

再证唯一性。如果 $x_2$ 在该开区间有至少两个零点，
由零点在内部无聚点，可选其中相邻的两个。
对调两解，刚证的存在性就会迫使 $x_1$ 在这两个点之间有零点，
与 $t_1,t_2$ 相邻矛盾。
\end{proof}


图~\ref{fig:ode-tikz-sturm-separation} 取 $x_1=\sin t,x_2=\cos t$；右图使用本章的相位约定 $(x,y)=\rho(\sin\theta,\cos\theta)$。
\begin{center}
\begin{minipage}{0.96\linewidth}
\centering
% TikZ mathematical illustration; compiled with the main document.
\begin{tikzpicture}[>=Stealth,font=\small]
\begin{scope}[x=.8cm,y=1.15cm]
 \draw[->,gray] (-.25,0)--(6.95,0) node[right]{$t$};
 \draw[->,gray] (0,-1.4)--(0,1.5) node[above]{$x$};
 \draw[GreenMain,thick,domain=0:6.2832,samples=100] plot (\x,{sin(\x r)});
 \draw[GoldMain,thick,domain=0:6.2832,samples=100] plot (\x,{cos(\x r)});
 \foreach \t/\lab in {0/0,1.5708/{\pi/2},3.1416/\pi,4.7124/{3\pi/2},6.2832/{2\pi}}{
  \draw (\t,.05)--(\t,-.05);\node[below,font=\scriptsize] at (\t,-.1){$\lab$};
 }
 \foreach \t in {0,3.1416,6.2832}{\fill[GreenMain] (\t,0) circle(2pt);}
 \foreach \t in {1.5708,4.7124}{\fill[GoldMain] (\t,0) circle(2pt);}
 \node[GreenMain] at (2,1.45){$\sin t$};\node[GoldMain] at (4.4,1.45){$\cos t$};
\end{scope}
\begin{scope}[shift={(8,0)},scale=1.15]
 \draw[->,gray] (-1.4,0)--(1.5,0) node[right]{$x$};
 \draw[->,gray] (0,-1.4)--(0,1.5) node[above]{$y=x'$};
 \draw[GreenMain,thick] (0,0) circle(1);
 \draw[GreenMain,->,thick] (70:1) arc[start angle=70,end angle=25,radius=1];
 \draw[GoldMain,->] (90:.48) arc[start angle=90,end angle=40,radius=.48];
 \node[GoldMain] at (.28,.65){$\theta$};
 \fill (0,1) circle(1.6pt) node[above left]{$\theta=0$};
 \fill (0,-1) circle(1.6pt) node[below left]{$\theta=\pi$};
\end{scope}
\end{tikzpicture}
\captionof{figure}{同一方程的独立解零点交错；右图中每经过一个零点，相位穿过一个 $k\pi$。}
\label{fig:ode-tikz-sturm-separation}
\end{minipage}
\end{center}

\section{比较定理与零点间距}
\begin{theorem}[Sturm 比较]
设 $Q,q$ 连续，$Q\geq q$。若 $x$ 解 $x''+qx=0$ 并以 $a<b$ 为相邻零点，
则每个非零解 $z$（$z''+Qz=0$）在 $[a,b]$ 至少有一个零点。
若 $Q>q$ 在 $[a,b]$ 上处处成立，则零点可取在 $(a,b)$。
\end{theorem}
\begin{proof}
将第一解的相位取为 $\alpha(a)=0,\alpha(b)=\pi$。
若第二解在某端点为零，闭区间版本已经成立。
否则改变其符号并选相位分支，使 $0<\beta(a)<\pi$。
设 $f_q(t,s)=\cos^2s+q(t)\sin^2s$，则
$\alpha'=f_q(t,\alpha),\beta'=f_Q(t,\beta)$。

令 $w=\beta-\alpha$，把差分精确写成
\[
 \begin{aligned}
 w'&=c(t)w+d(t),\\
 d(t)&=(Q(t)-q(t))\sin^2\alpha(t)\geq0,
\end{aligned}
\]
其中
\[
 c(t)=\displaystyle\int_0^1\partial_s f_Q(t,\alpha(t)+\tau w(t))\,d\tau
\]
为连续函数。常数变易给出
\[
 w(t)=e^{\int_a^tc}\left[w(a)+\int_a^te^{-\int_a^sc}d(s)\,ds\right].
\]
当 $w(a)>0$，方括号严格为正，因此 $\beta(b)>\pi$。
又 $\beta(a)<\pi$，故在开区间内取到 $\pi$，得到零点。
这证明了第二个解在两端均不为零时的结论，结合端点为零的情形，即得闭区间上的结论。

现在假设 $Q>q$。若第二解在左端点为零，取 $\beta(a)=0$，则 $w(a)=0$。
$\alpha$ 在 $a$ 后立即进入 $(0,\pi)$，因此 $d$ 在一个右邻域严格为正；
上式于是给出 $w(t)>0$ 对所有 $t>a$ 成立，特别是 $\beta(b)>\pi$。
所以除了左端零点以外，内部还有一次经过 $\pi$。
如果第二解仅在右端点为零，可从 $b$ 反向时间重复这一论证，
或把 $t$ 换成 $a+b-t$ 后应用左端情形，同样得到内部零点。
因此严格比较时零点可取在开区间。
\end{proof}
\begin{corollary}
若 $q\geq\omega^2>0$，则 $I$ 内相邻零点距离不超过 $\pi/\omega$；
当 $I$ 包含一条右半轴时，非零解还有任意大的零点。
若 $q\leq\Omega^2$，$\Omega>0$，则相邻零点距离不小于 $\pi/\Omega$。
若 $q\leq0$，每个非零解至多有一个零点。
\end{corollary}
\begin{proof}
常系数比较方程的零点间距分别是 $\pi/\omega$ 和 $\pi/\Omega$。
若存在超过相应界限的无零点区间，可在其内部放入比较解的两个相邻零点，
与比较定理矛盾。当 $I$ 包含右半轴时，把比较区间移到任意远处，
也得第一种情形的无限多个零点。
最后，若 $x$ 在 $a<b$ 为零，乘方程以 $x$ 并积分得
\[
 \int_a^b(x')^2\,\dd t=\int_a^bq(t)x(t)^2\,\dd t\leq0.
\]
所以 $x'=0$，由端点值知 $x=0$，与非零解矛盾。
\end{proof}
一般方程 $x''+p(t)x'+q(t)x=0$ 在 $p\in C^1$ 时，令
\[
 x=e^{-\frac12\displaystyle\int p}u,
\]
得到
$u''+(q-p'/2-p^2/4)u=0$。乘子从不为零，因此保持零点。

\section{Sturm--Liouville 问题}
\begin{definition}
给定连续实函数 $q$ 和 $l>0$，寻找
\[
 x''+(q(t)+\lambda)x=0,\qquad x(0)=x(l)=0
\]
的非零解，称为这里的 Sturm--Liouville 两点边值问题。
这样的 $\lambda$ 称特征值，相应非零解称特征函数。
\end{definition}
\begin{theorem}
特征值可以严格递增地排列为 $\lambda_1<\lambda_2<\cdots\to\infty$，
每个特征空间一维。第 $k$ 个特征函数在 $(0,l)$ 恰有 $k-1$ 个零点。
若 $m\leq q\leq M$，则
\[
 (k\pi/l)^2-M\leq\lambda_k\leq(k\pi/l)^2-m,
 \qquad \lambda_k/k^2\to(\pi/l)^2.
\]
\end{theorem}
\begin{proof}
\textbf{第一步：把边值问题化为一个相位方程。}
令 $x(t,\lambda)$ 解初值问题
\[
 x(0,\lambda)=0,\qquad x'(0,\lambda)=1.
\]
任何非零解若在左端为零，其左端导数必不为零，
除以该导数后便与 $x(t,\lambda)$ 具有相同初值，因此由唯一性重合。
所以每个特征空间一维，特征值条件正是 $x(l,\lambda)=0$。
取 $\theta(0,\lambda)=0$，这个条件等价于 $\theta(l,\lambda)=k\pi$，$k\geq1$。

\textbf{第二步：证明末端相位对参数严格增加。}
参数可微依赖定理给出 $w=\partial_\lambda x$，并由对方程求导得到
$w''+(q+\lambda)w=-x$、$w(0)=w'(0)=0$。
逐项计算为
\[
 (x'w-xw')'=x''w-xw''=x^2.
\]
零初值使 
\[
 x'w-xw'=\displaystyle\int_0^t x(s,\lambda)^2ds.
\]
又由相位角的微分公式，
\[
 \partial_\lambda\theta(l,\lambda)
 =\frac{x'(l,\lambda)w(l,\lambda)-x(l,\lambda)w'(l,\lambda)}
 {x(l,\lambda)^2+x'(l,\lambda)^2}>0.
\]
分母不为零，分子积分严格为正，因为 $x'(0,\lambda)=1$。

\textbf{第三步：确定相位的取值范围并编号。}
当 $\lambda\leq-\max q$ 时，$q+\lambda\leq0$。
解从零以正速度出发；在保持正值的阶段有 $x''=-(q+\lambda)x\geq0$，
所以导数不能降到零，解不可能再次为零。
因此末端相位位于 $(0,\pi)$。
另一方面，给定任意正整数 $N$，可取 $\lambda$ 足够大，使
$q+\lambda\geq\omega^2$ 且 $N\pi/\omega<l$。
比较定理使 $x$ 在 $[0,l]$ 有至少 $N$ 次零点穿越，故末端相位趋于无穷。
连续性与严格单调性说明每个水平 $k\pi$ 恰穿过一次，给出严格递增的 $\lambda_k$。
相邻零点角增为 $\pi$，故第 $k$ 个特征函数恰有 $k-1$ 个内部零点。
若 $\lambda_k$ 有有限聚点，末端相位在该紧参数区间有界，
但其值 $k\pi$ 无界，矛盾，所以 $\lambda_k\to\infty$。

\textbf{第四步：比较常系数方程以估计特征值。}
记 $\mu_k=(k\pi/l)^2$。
在 $\lambda=\mu_k-M$ 时有 $q+\lambda\leq\mu_k$，
由同初始角的标量比较，末端相位不超过常系数方程的 $k\pi$。
在 $\lambda=\mu_k-m$ 时有反向不等式，末端相位不小于 $k\pi$。
由末端相位严格随 $\lambda$ 增加，得到
$\mu_k-M\leq\lambda_k\leq\mu_k-m$。
两边除以 $k^2$ 后，常数误差趋于零，夹逼得到所述渐近式。
\end{proof}

\begin{definition}[平方平均空间与完全性]\label{def:l2-completion}
先在闭区间上的连续实函数空间 $C([0,l])$ 上定义
\[
 \langle f,g\rangle=\int_0^l f(t)g(t)\,dt,
 \qquad\|f\|_2=\langle f,f\rangle^{1/2}.
\]
将这个内积空间按 $\|\cdot\|_2$ 完备化，得到 $L^2(0,l)$。
具体地，它的元素可由连续函数的 Cauchy 列表示；两列等价，是指它们之差的 $2$ 范数趋于零。
内积定义为代表列内积的极限，Cauchy--Schwarz 不等式保证该极限存在且与代表列无关。
一个正交系的有限线性组合在 $2$ 范数下稠密时，称该正交系完全。
本节出现的特征函数都是经典光滑函数，因此其内积直接用上面的通常积分计算。
\end{definition}

\begin{proposition}
不同特征值的实特征函数在 $L^2(0,l)$ 中正交。
\end{proposition}
\begin{proof}
分别把两方程乘以另一个解后相减积分，得到
\[
 (\lambda_i-\lambda_j)\int_0^lx_i x_j\,\dd t
 =[x_i x_j'-x_i'x_j]_0^l=0.
\]
端点条件使边界项消失，而特征值不同，所以积分为零。
\end{proof}
\begin{definition}[权内积与形式自伴]
对连续正权 $w$，权内积为 $\langle f,g\rangle_w=\displaystyle\int_0^l wf g\,dt$。
表达式 $Ly=-w^{-1}[(py')'+qy]$ 配上零端点条件时，
两次分部积分给出 $\langle Lf,g\rangle_w=\langle f,Lg\rangle_w$；这称其在该边界条件下形式自伴。
这里“形式”指对光滑零端点函数的积分恒等式，不省略无穷维算子定义域的问题。
\end{definition}
方程 $(p x')'+(q+\lambda w)x=0$（$p\in C^1$，$q,w$ 连续且 $p,w>0$）
宜取相平面坐标 $(x,px')$；同样的 Wronskian 与相位比较可以使用，正交性采用上述权内积。
